% Chapter 4 — ResOpNet. Pure chapter content: no documentclass or preamble.
% Research proposal and mathematical specification, not a report of completed experiments.
% The same file is intended for the standalone article and the DHNTeam book.

\section{ResOpNet: Motivation and Research Objective}
\label{sec:resop-motivation}

Numerical integration can be viewed as two tasks: evaluating a simple discrete
sum and accounting for the signed area missed by that sum. The first task is
already solved efficiently by classical quadrature. The second is the object of
this chapter.

The proposed \emph{Residual Operator Network} (ResOpNet) separates these tasks:
\begin{equation}
    \boxed{\widehat{\mathcal I}_{\theta,Q}(f)
        = Q(f)+R_{\theta,Q}(S(f)).}
    \label{eq:resop-principle}
\end{equation}
Here $\mathcal I$ is the desired integration operator, $Q$ is an explicitly
chosen baseline quadrature operator, $S$ supplies observations of the input
function, and $R_{\theta,Q}$ predicts the signed difference between the true
integral and the baseline result.

The analogy with residual networks is architectural. Rather than asking a
network to relearn an inexpensive, known operation, we preserve that operation
as a separate path and train only its correction. ResNet uses an identity
shortcut; here the shortcut is a \emph{numerical operator}. The additive
decomposition by itself is not a claim of algorithmic novelty or improved
accuracy: the useful question is whether the learned correction exploits
information about a restricted family of functions that a general-purpose
quadrature method does not use.

Our primary practical objective is:
\begin{quote}
    At the same verified integration accuracy, reduce the number of expensive
    function evaluations and/or the total execution time relative to strong
    classical quadrature baselines.
\end{quote}
A second objective is to define reusable interfaces for function families,
local residual operators, sampling policies, and fair numerical benchmarks.
A larger research community or registry is an aspiration, not a result
established by this chapter.

\subsection{What is and is not being learned}

A composite rectangle or trapezoid sum is inexpensive once its samples are
available. We therefore distinguish \emph{cell-boundary samples}, used by the
classical quadrature path, from \emph{cell-interior samples}, given to the
learning path. The learning path may inspect many interior observations. There
is no requirement that it receive only a few values, or that every cell have
the same physical width.

In the minimal prototype, the baseline quadrature is fixed and exact as a
computation on its samples. An optional later mode replaces a costly baseline
operator by its own learned approximation.
% 这句话是不是不太需要
That optional mode is conceptually
separate from the first experiment; mimicking a cheap weighted sum with an MLP
is not in itself a speed optimization.

\section{Local Cells and the Reference Interval}
\label{sec:resop-reference}

Let $f:[a,b]\to\mathbb R$ and choose a mesh
\[
    a=x_0<x_1<\cdots<x_N=b,\qquad h_i=x_{i+1}-x_i>0.
\]
The exact contribution of the $i$th cell is
\[
    I_i(f)=\int_{x_i}^{x_{i+1}}f(x)\,dx.
\]
Every physical cell is pulled back to the same reference interval:
\begin{equation}
    A_i:[0,1]\longrightarrow[x_i,x_{i+1}],
    \quad t\longmapsto x_i+h_it,
    \qquad g_i(t):=f(A_i(t)).
    \label{eq:resop-affine}
\end{equation}
The change of variables gives the exact identity
\begin{equation}
    \boxed{I_i(f)=h_i\int_0^1 g_i(t)\,dt.}
    \label{eq:resop-cell-integral}
\end{equation}
Thus a single operator on $[0,1]$ can be reused for every cell. The factor
$h_i$ is restored only after evaluating the reference-cell operator.
Affine normalization fixes the domain, but it does \emph{not} make all
pulled-back functions statistically alike: frequencies, derivatives, and
singular features change as $h_i$ varies.

\subsection{Two sampling systems, with different responsibilities}
\label{sec:resop-two-samplings}

For clarity, use the composite trapezoid path first. It consumes the cell
endpoints
\begin{equation}
    E_i(f)=\bigl(f(x_i),f(x_{i+1})\bigr)
    =\bigl(g_i(0),g_i(1)\bigr).
\end{equation}
Choose interior reference locations
\[
    0<\tau_1<\cdots<\tau_m<1.
\]
The residual network consumes a separate interior observation vector
\begin{equation}
    S^{\mathrm{int}}_m(g_i)
    =\bigl(g_i(\tau_1),\ldots,g_i(\tau_m)\bigr)
    =\bigl(f(x_i+h_i\tau_j)\bigr)_{j=1}^{m}.
    \label{eq:resop-interior-sampling}
\end{equation}
These additional samples may be regular, nonuniform, multiscale, or selected
adaptively. The endpoint observations may also be supplied to the network
for detrending. The distinction concerns \emph{their computational role}, not
a rule that information must never be shared.

For the reference function $g$, define the normalized trapezoid operator
\begin{equation}
    Q_{\mathrm T}(g)=\frac{g(0)+g(1)}{2}.
    \label{eq:resop-trapezoid}
\end{equation}
The physical cell's classic contribution is $h_iQ_{\mathrm T}(g_i)$.
Other baselines, including rectangle, midpoint, Simpson, or Gaussian
quadrature, can later be specified by their nodes and weights on $[0,1]$.
When a baseline requires interior nodes, those evaluations are counted
explicitly in the comparison.

\section{The Exact Residual Operator}
\label{sec:resop-exact}

For a fixed reference quadrature rule $Q$, define the signed local residual
\begin{equation}
    \boxed{
        \mathcal R_Q(g)
        :=\int_0^1g(t)\,dt-Q(g).
    }
    \label{eq:resop-exact-residual}
\end{equation}
This is a map from an appropriate function space to $\mathbb R$; it is not
the residual of a PDE loss function, nor is it merely an asymptotic order
symbol. It is the \emph{actual signed quadrature error} of a given input.

If $Q$ is linear, then $\mathcal R_Q$ is linear on its domain. A generic
nonlinear neural network approximating $\mathcal R_Q$ is not automatically
linear.

A trained model
\[
    R_{\theta,Q}:\mathbb R^m\longrightarrow\mathbb R
\]
can be used when the sensor locations and function family are fixed. More
generally, it may accept sensor positions, endpoint values, scale, a method
descriptor, or family-specific features:
\[
    R_{\theta,Q}\bigl(S(g),\,\tau,\,h,\,\mathrm{context}\bigr).
\]
The fundamental physical-cell prediction is
\begin{equation}
    \boxed{
        \widehat I_i(f)
        =h_i\left[Q(g_i)+R_{\theta,Q}(S(g_i))\right].
    }
    \label{eq:resop-local-prediction}
\end{equation}
The full integral and all mesh-node antiderivatives are then obtained by
ordinary accumulation:
\begin{align}
    \widehat I(f)   & =\sum_{i=0}^{N-1}\widehat I_i(f), \\
    \widehat F(x_0) & =0,\qquad
    \widehat F(x_k)=\sum_{i=0}^{k-1}\widehat I_i(f)
    \quad (1\leq k\leq N).
    \label{eq:resop-prefix}
\end{align}
The prefix sum is \emph{not} another learned operator. Estimating values
between mesh nodes requires a separately specified within-cell reconstruction.

\subsection{Detrending: give the network only the missing curve}

The trapezoid rule integrates every affine function exactly. Let
\begin{equation}
    \ell_g(t)=(1-t)g(0)+tg(1),
    \qquad d_g(t):=g(t)-\ell_g(t).
    \label{eq:resop-detrend}
\end{equation}
Since $\int_0^1\ell_g(t)\,dt=Q_{\mathrm T}(g)$,
\begin{equation}
    \boxed{
        \mathcal R_{\mathrm T}(g)=\int_0^1d_g(t)\,dt,
        \qquad d_g(0)=d_g(1)=0.
    }
    \label{eq:resop-detrended-residual}
\end{equation}
Thus the learned task is literally the signed area between the curve and
its endpoint secant line.

For fixed interior positions, define
\[
    z(g):=\bigl(d_g(\tau_1),\ldots,d_g(\tau_m)\bigr)\in\mathbb R^m.
\]
A simple exactness-preserving realization is
\begin{equation}
    R_{\theta,\mathrm T}(z)=N_\theta(z)-N_\theta(0).
    \label{eq:resop-zero-constraint}
\end{equation}
It gives $R_{\theta,\mathrm T}(0)=0$ exactly, so the corrected formula
remains exact on affine functions regardless of training. Furthermore,
adding any affine function to $g$ leaves $z(g)$ unchanged; the baseline
integrates that added part exactly. This construction does \emph{not}
guarantee the linearity of $R_\theta$ on general inputs.

\subsection{A minimal schematic}

The intended computation is:
\[
    \begin{array}{c}
        f\ \longrightarrow\ \{[x_i,x_{i+1}]\}_{i=0}^{N-1}
        \\[2mm]
        \Big\downarrow\ \text{pull each cell back to }[0,1]
        \\[2mm]
        g_i
        \quad\longrightarrow\quad
        \begin{cases}
            \text{endpoints}\longrightarrow Q(g_i), \\
            \text{interior samples}\longrightarrow
            R_{\theta,Q}(S(g_i)).
        \end{cases}
        \\[2mm]
        \Big\downarrow\quad h_i\bigl(Q+R_{\theta,Q}\bigr)
        \\[2mm]
        \widehat I_i
        \quad\longrightarrow\quad
        \text{prefix sums and integral values}.
    \end{array}
\]
The baseline and residual branches share a reference function but have
distinct jobs. This decomposition is the design principle to preserve in
all implementations.

\section{What Classical Error Analysis Tells Us}
\label{sec:resop-regularity}

The notation $O(h^p)$ is a bound on how errors scale, not a computation of
the actual signed error for a specific cell. ResOpNet attempts to predict
that signed error from information about a restricted function family.
Classical error analysis remains useful for choosing the target, its scale,
and the necessary safeguards.

\subsection{Rectangle rule on a $C^1$ function}

For $f\in C^1([x,x+h])$, let
\[
    Q_h^{\mathrm L}(f;x)=hf(x),\qquad
    \mathcal R_h^{\mathrm L}(f;x)
    =\int_x^{x+h}f(s)\,ds-hf(x).
\]
Integration of $f'$ and interchange of integrals give
\begin{equation}
    \mathcal R_h^{\mathrm L}(f;x)
    =h^2\int_0^1(1-t)f'(x+ht)\,dt.
    \label{eq:resop-rectangle-error}
\end{equation}
Consequently,
\[
    \mathcal R_h^{\mathrm L}(f;x)
    =\frac{h^2}{2}f'(x)+o(h^2)
\]
as $h\to0$, with uniformity on a fixed compact domain for fixed $f\in C^1$.
This suggests learning a suitably normalized local residual rather than
a target whose magnitude vanishes rapidly with cell width.

\subsection{Trapezoid rule on $C^1$ and $C^2$ functions}

For the physical trapezoid contribution $Q_h^{\mathrm T}$, one obtains,
for $f\in C^1$,
\begin{equation}
    \mathcal R_h^{\mathrm T}(f;x)
    =h^2\int_0^1\left(\frac12-t\right)f'(x+ht)\,dt.
    \label{eq:resop-trap-c1}
\end{equation}
Writing
\[
    \omega_{f'}(h)
    =\sup_{\substack{s,t\in[a,b]\\|s-t|\leq h}}
    |f'(s)-f'(t)|
\]
gives
\begin{equation}
    |\mathcal R_h^{\mathrm T}(f;x)|
    \leq \frac{h^2}{4}\,\omega_{f'}(h).
    \label{eq:resop-trap-modulus}
\end{equation}
For a fixed continuously differentiable function this is $o(h^2)$
locally, but the full unbounded class $C^1$ has no common modulus of
continuity or a universal higher convergence rate.

For $f\in C^2$, the Peano-kernel representation is
\begin{equation}
    \mathcal R_h^{\mathrm T}(f;x)
    =-\frac{h^3}{2}
    \int_0^1t(1-t)f''(x+ht)\,dt,
    \label{eq:resop-trap-c2}
\end{equation}
and hence
\begin{equation}
    |\mathcal R_h^{\mathrm T}(f;x)|
    \leq\frac{h^3}{12}\|f''\|_{\infty,[x,x+h]}.
    \label{eq:resop-trap-bound}
\end{equation}
For example, $f(x)=x^2$ on $[0,h]$ gives
\[
    \int_0^h x^2\,dx-\frac h2\bigl(f(0)+f(h)\bigr)
    =-\frac{h^3}{6}.
\]
The network's target is the number $-h^3/6$, not merely the statement
that the error is $O(h^3)$.

\subsection{Information limits and honest guarantees}

Suppose a network sees only finitely many point evaluations. Even among
smooth functions, distinct functions can agree at every observed point
while having different integrals. Therefore no model, whatever its
capacity, can recover every integral over unrestricted $C^1$ or $C^2$
inputs exactly from those samples.

To make the task learnable we must specify a \emph{function family},
such as a bounded, parameterized family or a distribution of functions.
Accuracy on that family can be investigated experimentally; a worst-case
guarantee requires additional regularity and magnitude assumptions.
Affine normalization alone does not remove this information limit.

\section{Training the Reference-Cell Residual Operator}
\label{sec:resop-training}

\subsection{Generating supervised targets}

Choose a function family $\mathcal F$ and a distribution of its parameters.
For a training function $f$ and randomly selected cell $[x,x+h]$, form
$g(t)=f(x+ht)$. Generate the endpoint and interior observations separately.
Compute a trusted reference integral $I_{\mathrm{ref}}(g)$, analytically
when possible or by independently verified high-accuracy quadrature.
The signed reference target is
\begin{equation}
    r_Q(g)=I_{\mathrm{ref}}(g)-Q(g).
    \label{eq:resop-label}
\end{equation}
If a detrended trapezoid architecture is used, feed $z(g)$ from
\eqref{eq:resop-detrend}; otherwise feed a documented observation vector.

The elementary supervised objective is
\begin{equation}
    \boxed{
        \mathcal L_{\mathrm{res}}(\theta)
        =\frac1M\sum_{j=1}^{M}
        \left|R_{\theta,Q}(S(g_j))-r_Q(g_j)\right|^2.
    }
    \label{eq:resop-loss}
\end{equation}
Equivalently one may minimize
$|Q(g)+R_{\theta,Q}(S(g))-I_{\mathrm{ref}}(g)|^2$ directly.
Because $Q$ is a fixed path, the gradient updates the residual network,
not the known quadrature formula.

For variable physical widths, training on reference-cell residuals keeps
the explicit Jacobian $h$ outside the network. If family parameters,
frequencies, or other information vary with $h$, conditioning on that
context may still be necessary. Scale generalization is an experiment
to be measured, not assumed.

\subsection{Training splits and contamination control}

Split by \emph{complete generated functions and underlying parameters},
not by randomly shuffling cells cut from one function. Keep independent
tests for (i) new functions from the same family, (ii) new ranges of $h$,
(iii) new parameter regimes, and (iv) new families when transfer is claimed.
Record seeds, sampling locations, reference tolerances, software versions,
and the exact model checkpoint.

The cost of expensive reference integration belongs to offline training
and dataset generation. Deployment comparisons should report both the
per-query advantage and, when relevant, the number of future queries
needed to amortize training and data-generation costs.

\subsection{Residual correction versus error certification}

The model predicts $\widehat r$, while the remaining error is
\[
    e(g)=\int_0^1g(t)\,dt-Q(g)-\widehat r(g).
\]
These are different quantities. An empirical uncertainty score, a second
model, or agreement between coarse and refined meshes may help identify
difficult cases, but is not by itself a rigorous error certificate.
If a tolerance is advertised, specify whether it is a measured statistical
target or a mathematically certified bound. Unsupported inputs can be
handled by an explicitly implemented classical fallback.

\section{Operator Models and Replaceable Backbones}
\label{sec:resop-zoo}

The common object is the \emph{local signed residual operator}, not a
specific neural-network architecture. The simplest instance is an MLP
\[
    R_\theta:\mathbb R^m\to\mathbb R,
\]
using detrended interior samples. Possible alternatives are:
\begin{itemize}
    \item A compact MLP or convolutional residual corrector on fixed nodes.
    \item A DeepONet-style model when observations and query coordinates
          vary, or when a residual \emph{function} rather than one scalar
          per cell is required.
    \item An FNO-based model for functions sampled on regular grids and
          families with exploitable spectral structure. The number of
          retained modes and frequency range must be tested explicitly.
    \item Wavelet or multiscale encoders for localized features; specialized
          phase-aware methods for highly oscillatory families.
    \item Adaptive samplers that request extra evaluations in the cells
          where expected gains justify their costs.
\end{itemize}
For a scalar target on a handful of nodes, a large operator network may
be unnecessary. Architecture comparisons should focus on accuracy versus
\emph{total} cost rather than sophistication.

\subsection{Two distinct modes for the baseline path}

The initial, preferred implementation is
\begin{equation}
    \widehat{\mathcal I}=Q_{\mathrm{classical}}+R_\theta.
    \label{eq:resop-classical-mode}
\end{equation}
One may later consider
\begin{equation}
    \widehat{\mathcal I}=Q_\phi+R_{\theta,\phi},
    \label{eq:resop-learned-mode}
\end{equation}
where $Q_\phi$ is a separately learned approximation to an expensive
baseline. In the second mode $R_{\theta,\phi}$ must correct the error of
$Q_\phi$, \emph{not} merely the error of the original $Q$.

Learning a simple trapezoid weighted sum is unlikely to beat its direct
evaluation if the same function samples are already available. The
learned-baseline mode makes sense only if it reduces expensive evaluations
or substitutes for a genuinely costly operation.

\section{Function Families, Sampling Policies, and Adaptivity}
\label{sec:resop-families}

Different families may deserve different sensors and encoders while
sharing the same mathematical contract.

\subsection{First candidate families}

\begin{enumerate}
    \item \textbf{Smooth functions:}
          $f(t)=a+bt+ct^2+d\sin(2\pi\omega t+\varphi)$, with explicitly bounded
          parameter ranges. Use this to debug targets, verify affine exactness,
          and establish strong classical baselines.
    \item \textbf{Continuously differentiable but less smooth functions:}
          $f(t)=a+bt+c|t-\xi|^{3/2}$, $0<\xi<1$, as a simple $C^1$ example
          whose second derivative is singular at $\xi$. This separates the role
          of regularity from a purely smooth training set.
    \item \textbf{Oscillatory families:}
          $f(t)=A(t)\sin(\omega t+\varphi)$ with a controlled envelope $A$,
          frequency interval, and phase distribution. Compare local MLP,
          Fourier-aware, and phase-aware correctors.
    \item \textbf{Localized features:}
          narrow smooth bumps or steep transitions, with their width and
          location recorded. These test adaptive and nonuniform sampling.
\end{enumerate}
These are \emph{parameterized subsets} of regularity classes, not all of
$C^1$ or $C^2$. Fractal or highly irregular families are later research
subjects and will require their own definitions, information budgets,
and diagnostics.

\subsection{Nonuniform and adaptive cells}

For variable cell widths the affine rule \eqref{eq:resop-affine} is unchanged.
A sampling policy may choose internal $\tau_j$ conditional on prior
observations. Such extra function evaluations are part of the cost.

A useful future system combines (i) a signed correction estimate with
(ii) a separate prediction or check of the \emph{remaining} corrected
error. Only the second quantity should guide refinement decisions.
If it is unreliable, a conservative fallback and independent refined
quadrature may be preferable. The framework must not claim that a
large cell is safe simply because its residual network returns a small
number.

\section{Benchmark Protocol: Can It Beat Classical Quadrature?}
\label{sec:resop-benchmark}

The main research hypothesis is empirical: after offline learning,
a reusable local corrector might exploit family-specific information so
that coarser meshes or fewer expensive evaluations reach an accuracy that
otherwise needs more work. This should be demonstrated against strong
existing methods rather than inferred from training loss.

\subsection{Baselines and controlled comparisons}

For each family and budget, compare:
\begin{enumerate}
    \item Rectangle, trapezoid, midpoint, Simpson, and suitable
          Gauss-type rules.
    \item Adaptive classical integration at the same requested tolerance.
    \item A direct network that predicts the cell integral without a
          quadrature shortcut, using matched training data and capacity.
    \item The proposed $Q+R_\theta$, with and without detrending.
    \item Alternative residual architectures and sampling policies when
          their additional complexity is justified.
\end{enumerate}
A learned correction may rediscover a familiar higher-order rule.
For example, on three equally spaced nodes with spacing $h$,
\[
    Q_{\mathrm S}-Q_{\mathrm T}
    =\frac h6(-f_0+2f_1-f_2).
\]
Beating trapezoid through this identity alone would not demonstrate
an advantage over Simpson's formula.

\subsection{Metrics and cost accounting}

Report the final signed error, absolute and relative error (with care
when the true integral is near zero), high error quantiles, and the error
distribution on complete unseen functions. Also report:
\begin{itemize}
    \item Number of distinct evaluations of $f$, including every
          internal or adaptively selected sample.
    \item CPU/GPU inference time, classical accumulation time,
          memory use, and end-to-end wall-clock time.
    \item Accuracy at matched function-evaluation budgets, as well as
          time required to reach a matched \emph{verified} accuracy.
    \item Error versus cell width and distribution shift, including
          explicit failure cases.
    \item Training-data and training-compute costs separately from
          repeated-query inference.
\end{itemize}
Fast inference does not imply fast integration if the extra internal
samples dominate the runtime. Conversely, a small learned corrector may
be valuable for expensive functions, repeated queries, or challenging
special-purpose families.

If the predicted \emph{physical-cell} residual has error
$|R_i-\widehat R_i|\leq C h_i^{p+1}$ uniformly, then
\[
    \left|\sum_i(R_i-\widehat R_i)\right|
    \leq C\sum_i h_i^{p+1}
    \leq C(b-a)h_{\max}^{p}
\]
for $p>0$. This explains why predicting a signed correction with
controlled \emph{remaining} error matters. Such a bound requires an
actual justified hypothesis on the predictor; it cannot be deduced from
a small training loss alone.

\section{An Open ResOpNet Research Ecosystem}
\label{sec:resop-ecosystem}

Beyond the first model, the intended project is an open protocol that
allows researchers to introduce new function families, sampling methods,
and residual operators without replacing the shared evaluation tools.
Three components are proposed.

\subsection{Family Registry}

Each family entry specifies its domain, generator or dataset, parameter
distribution, regularity assumptions, supported scales, reference-value
method, and train/validation/test splits. A family may be deterministic,
stochastic, smooth, oscillatory, singular, or domain-specific; the entry
must identify its actual limits rather than claim universality.

\subsection{Operator Zoo}

Each submission specifies its input observation schema, sampling policy,
baseline quadrature $Q$, learned correction $R_\theta$, parameter count,
training recipe, model weights, validity range, and associated paper.
An MLP, FNO, DeepONet, or non-neural residual approximation can all
participate if they satisfy the same input-output and evaluation contract.
Author attribution and model licensing remain attached to each entry.

\subsection{Accuracy and Speed Arena}

A public evaluation harness measures final integral error and full
computational cost on held-out test functions. Separate leaderboards may
track smooth, oscillatory, low-regularity, costly-evaluation, and adaptive
settings. Submission rules should prevent leakage from hidden tests,
declare hardware and evaluation budgets, and preserve previous benchmark
versions so reported results remain reproducible.

A minimal software contract could be expressed as:
\begin{verbatim}
Family.sample(seed) -> function_spec
Family.reference(function_spec, interval) -> value
Sampler.observe(function_spec, cell) -> (positions, values)
Baseline.integrate(cell, endpoint_or_rule_samples) -> value
Residual.predict(reference_positions, reference_values,
                 baseline_context) -> signed_reference_correction
Benchmark.evaluate(family, sampler, baseline, residual) -> report
\end{verbatim}
This is an interface proposal, not an implemented API. It deliberately
separates the sampling budget from the residual model to make comparisons
auditable.

Framework releases should be versioned and citable; algorithm-specific
contributions should receive independent credit. A shared standard is
valuable only if others can reproduce results and choose not to adopt
the proposed network without losing access to the benchmark.

\section{First Proof-of-Concept and Milestones}
\label{sec:resop-prototype}

The first experiment should remain small enough to understand completely.
It does not require scattering solvers, complex-valued networks, or a
full research platform.

\subsection{Minimal experiment on $[0,1]$}

\begin{enumerate}
    \item Implement a family generator with known analytic integrals
          (initially smooth polynomials and simple trigonometric mixtures).
    \item Divide each function into physical cells and perform the affine
          pullback \eqref{eq:resop-affine}.
    \item Compute the trapezoid baseline from endpoints only; obtain
          separate interior observations for a small MLP.
    \item Train on the \emph{signed detrended residual} from
          \eqref{eq:resop-detrended-residual}, with independent reference
          values as supervision.
    \item Produce plots of per-cell signed residuals, corrected integrals,
          prefix sums, error versus mesh size, and wall time versus
          achieved error.
    \item Compare against Simpson, Gauss-type, and adaptive baselines at
          explicitly matched total sample budgets; document both
          improvements and failures.
\end{enumerate}

\subsection{Acceptance criteria}

The initial result is a \emph{reproducible experiment}, not a claim that
the method beats classical quadrature. It is acceptable when the full
training and testing pipeline runs deterministically, labels and
summations are checked against analytic cases, sampling costs are
reported, and failures are visible.

A stronger subsequent claim requires consistent gains on held-out
instances at matched accuracy, including inference and function
evaluation costs. Only after that should the project expand to
more families, FNO/DeepONet alternatives, adaptive meshes, a model
registry, and an external challenge.

\section{Relation to Acoustic Scattering and Future Questions}
\label{sec:resop-outlook}

Numerical integration is the starting point because its classical
baseline and its local signed residual are directly observable with
trusted references. In a later stage, structured residual learning may
be useful in boundary integral equations for scattering, where kernel
oscillation, geometric features, and near-singular quadrature can make
individual evaluations expensive. Singular terms must be treated
explicitly; a corrector trained on smooth reference-cell functions is
not automatically transferable to a singular Helmholtz kernel.

The larger research questions are:
\begin{itemize}
    \item Which families admit accurate local correction from a fixed
          sampling budget, and which do not?
    \item Which sampling policies and operator backbones provide the
          best accuracy--time tradeoff for each family?
    \item How much can one model generalize over affine cell widths and
          outside its parameter distribution?
    \item Can the corrected result be accompanied by a reliable
          error estimate or a rigorous certification procedure?
    \item Does a common registry and benchmark lead to reusable methods
          beyond the first proof-of-concept?
\end{itemize}

\paragraph{Status.}
This chapter specifies the ResOpNet idea, its mathematical targets, and
its proposed experimental and community infrastructure. No speedup,
accuracy improvement, universal function-class guarantee, or benchmark
result is asserted prior to numerical validation.




\section{Reference}
\bibliographystyle{plain}  % or another style like unsrt, IEEEtran, etc.
\bibliography{refer}  % references.bib is the file name